\section{Distribution estimates for the sieve coefficients}
\label{sec:current-distribution}

\begin{definition}
A finite prefix policy specifies finitely many regions in the
coordinates $t_i=\log_Np_i$ of ordered prime prefixes and a sieve level
on each region. The decision at a prefix is fixed before its children
are refined. A selected branch is admissible when its full-square and
selected-complement conditions in
\eqref{coef:full-squares}--\eqref{coef:selected-squares} have the required
strict margins; an inherited branch keeps its parent decision.
A marked family is a nonnegative finite weighted family of prefixes
$B$ at a common total level $Q$, with residual level $Q/B$ and the
original prime ceiling retained. A marked query is a sifted count for
one such family. References below to an admitted coefficient mean a
coefficient obtained from these policies and finite marked transitions,
with the three-factor representations of
Appendix~\ref{app:coefficients}.
\end{definition}

\subsection{A local estimate with a moving residue}

The distribution problem is local in the complete product. We denote
the convolution lengths by $U,V$, write $x=UV$, and use $\mathfrak a$
for the residue; $N$ remains the global even target.
The estimate proved in Appendix~\ref{app:medium} is the following.

\begin{proposition}[Restricted medium convolution]
\label{prop:current-medium}
Fix $\vartheta=7/32$, $0<\sigma<1/100$, $c_0>0$, a divisor order
$B$, and $0\le\kappa\le\sigma/100$. Let $U,V,Q_1,Q_2,Q_3\ge1$. Suppose
\begin{align}
 c_0x&\le |\mathfrak a|\le x^{1+\kappa},&
 Q_1&<Ux^{-\sigma}, \notag\\
 U^2Q_2Q_3^2&<x^{1-8\sigma},&
 |\mathfrak a|^{\vartheta}UQ_1Q_2^5Q_3^2&<x^{2-15\sigma},
       \label{eq:current-medium-ranges}\\
 U^{1+2\vartheta}Q_1Q_2^{5-2\vartheta}Q_3^{2+2\vartheta}
      &<x^{2-15\sigma}.&&\notag
\end{align}
Let $\alpha$ be one-bounded, supported on an individual interval in
$[U,2U]$ and on
$P^-(n)>x^{1/(\log\log x)^3}$.  Assume its auxiliary-coprimality
Siegel--Walfisz estimate: for every $A>0$, every $k,q\ge1$ and
$(b_0,q)=1$,
\begin{equation}
 \left|\sum_{\substack{n\equiv b_0\pmod q\\(n,k)=1}}\alpha_n
 -\frac1{\varphi(q)}\sum_{(n,kq)=1}\alpha_n\right|
 \le C_AU\tau(k)^{C_0}(\log x)^{-A},
 \label{eq:current-aux-sw}
\end{equation}
with fixed $C_0$ and uniform family constants.  Let
$|\beta_m|\le\tau(m)^B$, $m\in[V,2V]$, and let each $\gamma_i$ be a
separate one-bounded sequence on $q_i\asymp Q_i$, supported on
$(q_i,\mathfrak a)=1$.  For
\[
 \Delta_{\ell}=\sum_{n,m}\alpha_n\beta_m
 \left({\one}_{nm\equiv\mathfrak a\pmod \ell}
       -\frac{{\one}_{(nm,\ell)=1}}{\varphi(\ell)}\right),
\]
the stated classical inputs of Appendix~\ref{app:classical} give
\begin{equation}
 \sum_{q_1,q_2,q_3}\gamma_1(q_1)\gamma_2(q_2)\gamma_3(q_3)
                    \Delta_{q_1q_2q_3}
       =O_A\bigl(x(\log x)^{-A}\bigr).
 \label{eq:current-medium-conclusion}
\end{equation}
The constant is uniform in the residue and families satisfying these
conditions.  Fixed endpoint multiples are absorbed by the strict
power margins.
\end{proposition}

\paragraph{Application of the local estimate.}
Fix one finite admitted coefficient family. Appendix~\ref{app:coefficients}
uses its full and selected-complement square conditions to construct three
separate modulus factors on whole atom boxes. Appendix~\ref{app:prime-products}
verifies the remaining hypotheses of Proposition~\ref{prop:current-medium}
and applies it to the finite prime-product families. The
\hyperlink{coef:medium-checks}{hypothesis checks in Appendix~\ref{app:coefficients}}
identify these steps. Low-rank pieces and the branches without a medium slot
have their separate estimates. The resulting common input for the sieve
constructions is the prime-product distribution estimate in the next subsection.

The three modulus coefficients in
Proposition~\ref{prop:current-medium} remain separate, and the estimate
allows their signs. In the application the medium slot is a rough
$1$ or M\"obius sequence on an individual chart, with a permitted
Mellin twist. Appendix~\ref{app:prime-products} proves the auxiliary
Siegel--Walfisz condition~\eqref{eq:current-aux-sw} for these sequences.

When the residue has size comparable to $x$, reciprocity leaves an
oscillatory phase. Appendix~\ref{app:medium} keeps this phase and bounds
its variation within the strict power margins
of~\eqref{eq:current-medium-ranges}. The spectral estimates, including
their Fourier normalizations, are stated in
Appendix~\ref{app:classical}.

\subsection{From the local estimate to prime products}

For a fixed finite smooth weight $W_N$ on $k$ labelled original primes,
all at least $cN^{a_*}$ for fixed $a_*>0$, put
\[
 X_W^{(q)}=\sum_{(p_1\cdots p_k,q)=1}W_N(p_1,\ldots,p_k),
 \qquad X_W=\sum W_N(p_1,\ldots,p_k).
\]
The product is $O(N)$.  The allowed geometric families have fixed
polylogarithmic mixed Euler-derivative bounds and no additional
arithmetic filter.  The conclusion used throughout the paper is
\begin{equation}
 \sum_q\rho_N(q)
 \left\{\sum_{p_1\cdots p_k\equiv N\pmod q}W_N(p_1,\ldots,p_k)
            -\frac{X_W^{(q)}}{\varphi(q)}\right\}
       =O_{A,W}\bigl(N(\log N)^{-A}\bigr).
 \label{eq:current-PP}
\end{equation}
The same signed coefficient $\rho_N$ occurs in both terms.  Its
admitted classes are the actual upper tree, the actual lower trees,
the marked queries and ordinary descendants of the finite $H/h$
construction, and the positive switched families on which the upper
coefficient is used.

The coefficient proof in Appendix~\ref{app:coefficients} first freezes
one finite causal policy.  The physical level $N^\lambda$ determines
its density and its residual $Q/B$.  A slightly larger auxiliary
support cover $N^{\ell_{\rm aux}}$ may be used to prove its
factorization; it does not change $\lambda$ or the integer coefficient.
For an ordered prime prefix write $t_i=\log_Np_i$ and
$S_j=\sum_{i=1}^jt_i$, and let $I$ be the subset of slots selected
for one factor of the modulus. The full-square conditions
\[
 S_{j-1}+2t_j\le\ell_{\rm aux}
\]
and every selected-complement square
\[
 \sum_{i\in I}t_i+
       \sum_{\substack{i<j\\i\notin I}}t_i+2t_j
       \le\ell_{\rm aux}\quad(j\notin I)
\]
are retained. After all selected slots have appeared, the second test
becomes the full-square test.  Descendant closure and the physical
prime ceiling therefore persist through the finite ordinary queries.

There is an exact tensor identity before estimating discrepancies.
For a common prime-label alphabet with multiplicities $m_I$, split
the slots into three groups and divide each group convolution by
$M_j=\prod_I m_{I,j}!$.  The resulting $\gamma_j$ is one-bounded.
With $M_{\rm all}=\prod_I m_I!$, the full normalized convolution is
\[
 \frac{M_1M_2M_3}{M_{\rm all}}
                       (\gamma_1*\gamma_2*\gamma_3).
\]
Marked multiplicities remain in the external scalar.  The small-prime
beta factor attaches by the unique small/rough factorization.  The
partition is fixed on a whole atom box before any arithmetic
discrepancy is inspected.  Global squarefreeness remains in the
physical coefficient; its removal only for analytic comparison has
the fixed-power error
$O((N+N^{\theta_*})N^{-\eta_0+o(1)})$, where $N^{\eta_0}$ is
the small-prime presieve cutoff and $\theta_*<3/5$ is the support
exponent of the finite coefficient construction.

Appendix~\ref{app:prime-products} makes a smooth complete-product trim
below $2N^{1-\kappa_0}$ and then an exact smooth dyadic partition.
The local scale is the complete product, with
\[
 cN^{1-\kappa_0}\le x\le CN,\qquad
 c_0x\le N\le x^{1+\kappa}.
\]
It is never a residual quotient.  On an original group at scale $R$,
the exact order-three Heath--Brown identity is obtained from
$\Lambda*(\delta-1*\mu_Z)^{*3}=0$ below $2R$, with
$Z=(2R)^{1/3}$:
\begin{equation}
 \Lambda=\sum_{j=1}^3(-1)^{j-1}\binom3j
                 \mu_Z^{*j}*1^{*(j-1)}*\log.
 \label{eq:current-HB}
\end{equation}
The smooth original chart is retained; no sharp coupled group boundary
is inserted.  The Fourier separation of the geometric profile is
charged by its integrable Fourier norm.  The actual rough slots have
the auxiliary SW estimate.  Medium branches use
Proposition~\ref{prop:current-medium}; low ranks use below-half
Bombieri--Vinogradov.  The one-long and two-long alternatives are the
direct Brun arguments of Appendix~\ref{app:brun}, keeping the exact
coprime mean and both Poisson zero-axis cancellations.  In particular,
negative lower Brun polynomials are not multiplied as though they
were positive.

Choose the fixed moment order $K_{\rm mom}$ before narrowing
\[
 0<\kappa\le\min\left(\kappa_*,\frac{\sigma}{200},
                           \frac1{2K_{\rm mom}}\right),\qquad
 0<\kappa_0<\frac{\kappa}{4(1+\kappa)}.
\]
Here $\kappa_*>0$ is fixed by the strict support margins in
Appendix~\ref{app:coefficients}. The complete local support has exponent $\Theta_{\rm loc}<3/5$;
the long-branch parameter $\zeta$ is selected with
$20\zeta+\Theta_{\rm loc}<1$ and
$40\zeta+3\Theta_{\rm loc}/2<1$.  All finite atom scalar masses and
the fixed logarithmic number of charts are paid by stronger local
logarithmic accuracy.  This proves \eqref{eq:current-PP} for the
admissible families and all sufficiently large even targets.  

\section{A positive upper sieve with its own density}
\label{sec:current-upper}

The upper construction uses
\[
 \rho_0=\frac{19}{96},\quad \mu_0=\frac{25}{128},\quad
 \Theta=\frac{299}{500},\quad \Lambda_{\rm cap}=\frac{153}{256},\quad
 v(t)=\min\left(\frac{1+t}{2},\frac{57}{64}-\frac{3t}{2}\right).
\]
Its causal cap changes are made on the actual prefix; every retained
ancestor decision is preserved.  Fourth, fifth and sixth changes are
successive replacements, whose raw density differences telescope.
Later eighth-rank changes occupy $30+15+13+29=87$ specified domains.
Same-stage domains have disjoint interiors; successive stages use
differences from their respective parent caps.  
Appendix~\ref{app:upper} supplies the cap formulas and separation proofs.

\subsection{Exit identity and literal residual densities}

For a finite allowed prime universe, let $\xi$ be multiplicative with
$\xi(1)=1$ and $0\le \xi(p)\le1$. The upper
tree, whose first exits have odd rank, satisfies the exact identity
\begin{equation}
 \sum_d\mu(d){\one}_{d\ {\rm retained}}\xi(d)
 =\prod_p(1-\xi(p))
   +\sum_{e\ {\rm odd\ first\ exit}}\xi(e)
                      \prod_{p<p_{\rm last}(e)}(1-\xi(p)).
 \label{eq:current-upper-exit}
\end{equation}
Partition the excluded divisors by their first excluded prefix; the
remaining smaller-prime subsets give the displayed Euler product.
Taking $\xi(p)={\one}_{p\mid n}$ proves a rough-indicator majorant for
every integer $n$, including zero and negative integers.

For the raw density, take
$g_N(p)={\one}_{p>2,\,p\nmid N}/(p-1)$ and
$V_N(y)=\prod_{2<p<y,\,p\nmid N}(1-1/(p-1))$.
At a fixed rank the normalized exit for $g_N$ lies between zero and
the corresponding exit for $g_0(p)={\one}_{p>2}/(p-1)$: a surviving
prefix acquires the multiplier
\[
 \prod_{\substack{p_{\rm last}\le p<y\\p\mid N}}(1-g_0(p))\le1.
\]
The preceding square tests bound the last prime of a rank-$j$ exit
below by $N^{\beta_{\min}/(3\cdot2^{j-1})}$ on positive compact
parameter sets.  Deleting target primes thus changes each fixed-rank
integral by $o(1)$ uniformly in $N$.  Positive exit-series convergence,
the one-sided bounds for the same raw Rosser weights, and then a
finite-rank sandwich yield
\begin{equation}
 \frac{V_N^+(R,y)}{V_N(y)}=F(s)+o(1)\quad(s>1),\qquad
 \frac{V_N^-(R,y)}{V_N(y)}=f(s)+o(1)\quad(s\ge2),
 \quad s=\frac{\log R}{\log y}.
 \label{eq:current-raw-density}
\end{equation}
The upper identity for $1<s\le3$ also follows directly by moving its
first cube to $R^{1/3}$.  This computes the density of the divisor coefficient itself.  The ordinary exit-series inputs and their normalization are
specified in Appendix~\ref{app:upper}.

Mertens gives, uniformly in the target,
\begin{equation}
 V_N(N^u)=\frac{2e^{-\gamma}C_N}{u\log N}(1+o(1)).
 \label{eq:current-mertens}
\end{equation}
For example the exact ratio to $2C_N\prod_{p<z}(1-1/p)$ is
\[
 \prod_{p\ge z}\left(1-\frac1{(p-1)^2}\right)^{-1}
 \prod_{\substack{p\mid N\\p\ge z}}\left(1-\frac1{p-1}\right),
\]
whose logarithm tends uniformly to zero at $z=N^u$, $u>0$ fixed.
With $\overline F=F/(2e^\gamma)$ and
$\overline f=f/(2e^\gamma)$, this is the normalization producing
$4\overline F(s)/u$ and $4\overline f(s)/u$ in the density functionals.

For a selected eighth cell, if $S_8=t_1+\cdots+t_8$ and the parent is bounded above by $R_6$, its cap $M$ satisfies
\[
 1<\frac{\lambda_{6,\rm parent}-S_8}{t_8}
  \le\frac{R_6-S_8}{t_8}<\frac{M-S_8}{t_8}<3.
\]
The same residual identity therefore gives the lower saving
\begin{equation}
 \frac{4(M-R_6)}{(R_6-S_8)(M-S_8)\prod_{i=1}^8t_i}.
 \label{eq:current-eighth-saving}
\end{equation}
Using the larger parent $R_6$ is conservative.  This density has no
Chen quarter or ordering factorial.  The selected domains have
$t_8>53/13824$, and all full and selected-complement squares remain
in force.

\subsection{Finite realization and density errors}

The ideal policy is approximated on a finite causal partition before
$N$ grows.  Small-last-prime tails tend to zero; proper affine walls
have vanishing mass, whereas open inherit/equality regions are
retained.  At even rough rank take outer endpoint tests and at odd
rank inner tests, preserving ancestry.  This produces a coefficient
$a_N(d)$ above the raw upper coefficient.  Attach the small-prime
beta upper polynomial to positive $a_N(d)$ and the beta lower
polynomial to negative $a_N(d)$.  If $B^-\le I_w\le B^+$, then
$a_N(d)B_{\rm sign}(n)\ge a_N(d)I_w(n)$ term by term.
The resulting physical $\rho_N^{\rm up}$ has unique small/rough decomposition,
squarefree support $q\le N^{\theta_*}$, a uniform bound
$|\rho_N^{\rm up}(q)|\le M$ for fixed $M$, $(q,N)=1$, and
\begin{equation}
 R_N(n):=\sum_{q\mid n}\rho_N^{\rm up}(q)
   \ge {\one}_{(n,P_N(N^{\rho_0}))=1}\ge0,
 \qquad P_N(y)=\prod_{\substack{p<y\\p\nmid N}}p.
 \label{eq:current-physical-upper}
\end{equation}
The finite sum also defines $R_N(0)$.

Restricting \eqref{eq:current-upper-exit} to primes above the presieve
gives the upper projection inequality
\[
 V_N(w)\sum_{d\ {\rm rough}}\lambda_N(d)g_N(d)
       \le\sum_{d\ {\rm full}}\lambda_N(d)g_N(d).
\]
Together with the absolute beta and box-error estimates this controls
the density of the same $\rho_N^{\rm up}$ from above.  Support inflation can
be kept below $1/1000$, giving $\theta_*=599/1000$ before the small
complete-scale adjustment.  The auxiliary support cover used in
Section~\ref{sec:current-distribution} does not change this density.

Let $E_{\rm up}$ be the raw upper-density endpoint and let $u$ be each of the
two allowances for finite approximation and coefficient realization.
Their exact rational values are given in
\eqref{eq:current-E}--\eqref{eq:current-u}. The directed estimates in
Appendix~\ref{app:numeric} give
\begin{equation}
                     K-E_{\rm up}-2u>9.6522\cdot10^{-8}>0.
                     \label{eq:current-upper-slack}
\end{equation}
Thus the physical coefficient in \eqref{eq:current-physical-upper} satisfies
\begin{equation}
 S_N:=\sum_q\frac{\rho_N^{\rm up}(q)}{\varphi(q)}
       \le(K+o(1))\frac{C_N}{\log N},
 \label{eq:current-upper-density}
\end{equation}
and it is precisely the coefficient to which
\eqref{eq:current-PP} applies.

\begin{proof}[Proof of Theorem~\ref{thm:positive-coefficient}]
Use the finite coefficient constructed above.
Equation~\eqref{eq:current-physical-upper} gives its bounded squarefree
support and~\eqref{eq:core-positivity}, and
\eqref{eq:current-upper-density} gives its density bound.
The factorization and parameter choices in
Appendix~\ref{app:coefficients}, followed by the prime-product estimate
in Appendix~\ref{app:prime-products}, establish
\eqref{eq:current-PP} for this coefficient and each weight in the
fixed finite collection. Choose their common refinement before taking
$N$ to infinity. Taking maxima of the finitely many constants and
thresholds gives the asserted uniformity.
\end{proof}

\begin{proposition}[Positive product consequence]
\label{prop:current-positive-product}
Fix a finite nonnegative labelled prime-product family of bounded
rank and uniformly bounded weights, with every original prime at least $N^{a_*}$ for some fixed
$a_*>0$, product $O(N)$, and only the admitted fixed geometric
restrictions.  Suppose its unsifted weighted mass $X_W$ satisfies
\[
 \limsup_{\substack{N\to\infty\\2\mid N}}
             \frac{\log N}{N}X_W\le J<\infty.
\]
Let $W_N(m)$ denote the induced weight on the product $m$. Then
\[
 \limsup_{\substack{N\to\infty\\2\mid N}}\frac1{H_N}
       \sum_{\substack{m<N\\N-m\text{ prime}}}W_N(m)\le KJ.
\]
The error estimate is absolute also when $J=0$.
\end{proposition}

\begin{proof}
For a prime $\ell$, squarefree support and boundedness give
\[
 \sum_{q:\ell\mid q}\frac{|\rho_N^{\rm up}(q)|}{\varphi(q)}
       \ll\frac{\log N}{\ell-1}.
\]
Hence, by a union bound over the fixed number of original prime slots,
\begin{equation}
 \left|\sum_q\frac{\rho_N^{\rm up}(q)}{\varphi(q)}
                         (X_W^{(q)}-X_W)\right|
       \ll X_WN^{-a_*}\log N=o(H_N).
 \label{eq:current-coprime-correction}
\end{equation}
No factor $N^{\theta_*}$ is introduced.  Small prime remainders cost
$N^{\rho_0+o(1)}$.  The smooth complete-product trim costs
$N^{1-\kappa_0+o(1)}$ after the full modulus sum; its product is
strictly below $N$, so no $\tau(0)$ occurs.  Positive majorization
before distribution gives
\[
 \sum_{m<N,\,N-m\ {\rm prime}}W_N(m)
 \le\sum_m W_N(m)R_N(N-m)+o(H_N)
 \le X_WS_N+o(H_N).
\]
Outer approximations with $m\ge N$ are harmless because
$R_N(n)\ge0$ for every integer.  Now apply
\eqref{eq:current-upper-density}.  Fixed geometric boundaries are
handled by positive smooth outer approximation, with $N$ tending to
infinity before the approximation error tends to zero.  At each
precision the same $\rho_N^{\rm up}$ is used in the pointwise inequality,
density and distribution formula.  It need not be the same
coefficient at different precisions.
\end{proof}

\section{The full Chen weight and paired deletion}
\label{sec:current-chen-weight}

We use $a_{\rm s},h_{\rm s},d_{\rm s}$ for the fixed sieve cutoffs;
the parameter $a$ in Proposition $(1+a)$ retains its introductory meaning. Put
\[
 {a_{\rm s}}=\frac{417}{5000},\qquad {h_{\rm s}}=\frac{19}{125},\qquad
 {d_{\rm s}}=\frac{59}{250},\qquad \mu_0=\frac{25}{128}.
\]
For $n\le N$ define
$Q_d(y;n)=\one_{d\mid n}\one_{P^-(n/d)\ge y}$.
All prime prefixes below have distinct entries.  Suppressing the
argument $n$, the eleven terms are
\begingroup\allowdisplaybreaks[4]
\begin{align*}
 Y_1&=Q_1(N^{a_{\rm s}}),\\
 Y_2&=Q_1(N^{h_{\rm s}}),\\
 Y_3&=\sum_{N^{a_{\rm s}}\le p<N^{1/3}}Q_p(N^{a_{\rm s}}),\\
 Y_4&=\sum_{N^{a_{\rm s}}\le p<N^{d_{\rm s}}}Q_p(N^{a_{\rm s}}),\\
 Y_5&=\sum_{N^{a_{\rm s}}\le p_1<p_2<N^{h_{\rm s}}}Q_{p_1p_2}(N^{a_{\rm s}}),\\
 Y_6&=\sum_{N^{a_{\rm s}}\le p_1<N^{h_{\rm s}}\le p_2<N^{d_{\rm s}}}Q_{p_1p_2}(N^{a_{\rm s}}),\\
 Y_7&=\sum_{N^{d_{\rm s}}\le p_1<p_2<(N/p_1)^{1/2}}Q_{p_1p_2}(p_2),\\
 Y_8&=\sum_{N^{a_{\rm s}}\le p_1<N^{1/3}\le p_2<(N/p_1)^{1/2}}
                                          Q_{p_1p_2}(p_2),\\
 Y_9&=\sum_{N^{h_{\rm s}}\le p_1<N^{d_{\rm s}}\le p_2<(N/p_1)^{1/2}}
                   Q_{p_1p_2}\bigl((N/(p_1p_2))^{1/2}\bigr),\\
 Y_{10}&=\sum_{N^{a_{\rm s}}\le p_1<p_2<p_3<p_4<N^{h_{\rm s}}}Q_{p_1p_2p_3p_4}(p_2),\\
 Y_{11}&=\sum_{N^{a_{\rm s}}\le p_1<p_2<p_3<N^{h_{\rm s}}\le p_4<N^{d_{\rm s}}}
                                          Q_{p_1p_2p_3p_4}(p_2).
\end{align*}
\endgroup
In particular $Y_{11}$ has its full last range.  Set
\begin{equation}
 W=3Y_1+Y_2-Y_3-Y_4+Y_5+Y_6-2Y_7-Y_8-Y_9-Y_{10}-Y_{11}.
 \label{eq:current-full-weight}
\end{equation}

The terms $3Y_1+Y_2$ supply the rough-count baseline at the two
cutoffs. The one-prefix subtractions $Y_3,Y_4$ remove overcounting,
while $Y_5,Y_6$ restore the selected two-prefix contributions.
The terms $2Y_7,Y_8,Y_9$ subtract the remaining three-or-more-factor
patterns, and the four-prefix terms $Y_{10},Y_{11}$ compensate the
growth of the positive pair counts when several small factors occur.
Their combination gives the pointwise inequality below. The later
paired deletion links $Y_6$ and $Y_{11}$ by an injection, so that the
change preserves that inequality.

\begin{lemma}[Pointwise full-weight inequality]
\label{lem:current-full-weight}
For squarefree $n\le N$,
\[
                  W(n)\le4{\one}_{\Omega(n)\le2}.
\]
\end{lemma}

\begin{proof}
If $n$ is not $N^{a_{\rm s}}$-rough, all eleven terms vanish.  For $Y_9$, the
only less immediate case, its quotient cutoff exceeds
$N^{(1-{d_{\rm s}})/4}>N^{a_{\rm s}}$.  Suppose therefore that $n$ is squarefree and
$N^{a_{\rm s}}$-rough.  Let $r,s,t$ count its prime factors below $N^{h_{\rm s}},N^{d_{\rm s}},N^{1/3}$,
respectively.  Its first six counts are
\[
 (Y_1,Y_2,Y_3,Y_4,Y_5,Y_6)
      =\left(1,{\one}_{r=0},t,s,\binom r2,r(s-r)\right).
\]
If $r\ge2$, the quotient sifted at the second selected prime forces
the first two primes in $Y_{10}$ and $Y_{11}$ to be the two least
prime factors of $n$.  Thus
\[
 Y_{10}=\binom{r-2}{2},\qquad Y_{11}=(r-2)(s-r),
\]
with the binomial zero at $r=2,3$.  Substitution in the first six
weighted terms minus $Y_{10}+Y_{11}$ gives $s-t\le0$; the other
subtractions help.

For $k=\Omega(n)\ge3$, let $q_1<q_2<\cdots$ be the factors of $n$.
Then $q_1q_2^2<n\le N$, so $t\ge1$,
$Y_7={\one}_{s=0}$ and $Y_8={\one}_{t=1}$.
When $r\le1$, the remaining upper bounds for $W$ are
\[
\begin{array}{c|c}
 \text{case}&\text{upper bound}\\ \hline
 r=1&2-t-{\one}_{t=1}\\
 r=0,\ s\ge2&4-s-t\\
 r=s=0&2-t-{\one}_{t=1}\\
 r=0,\ s=1,\ t\ge3&3-t\\
 r=0,\ s=1,\ t\le2&3-t-1-{\one}_{t=1}.
\end{array}
\]
In the last case, $k\ge4$ would imply
$n\ge N^{{h_{\rm s}}+{d_{\rm s}}+2/3}>N$, so $k=3$; moreover
$q_1q_2q_3^2\ge N^{{h_{\rm s}}+{d_{\rm s}}+2/3}>N$.  The pair $(q_1,q_2)$ therefore
contributes to $Y_9$, giving the extra subtraction displayed there.
All the bounds are nonpositive.  For $k\le2$ direct substitution
gives $W\le4$; the unit has $W(1)=4$.
\end{proof}

Let $\mathscr E=\mathscr E_0\cup \mathscr E_1$ be the disjoint union of the rational rectangles
specified in Appendix~\ref{app:numeric}. The two families have $6091$ and
$11481$ members, respectively, in the coordinates
\[
 (t,u)=(\log_Np_2,\log_Np_1),\qquad {a_{\rm s}}\le u<{h_{\rm s}}\le t<{d_{\rm s}}.
\]
The rectangles have disjoint interiors; shared edges have zero limiting
mass and are assigned once in the finite partition.
Write $Y_6(\mathscr E)$
for its pair terms and $Y_{11}(\mathscr E)$ for the four-prefix terms whose
last two selected primes $(p_3,p_4)$ have
$(\log_Np_4,\log_Np_3)\in \mathscr E$.  The weight is
\begin{equation}
                         W_{\mathscr E}=W-Y_6(\mathscr E)+Y_{11}(\mathscr E).
 \label{eq:current-paired-weight}
\end{equation}

\begin{lemma}[Paired injection]
\label{lem:current-paired-injection}
For squarefree $n\le N$, $Y_{11}(\mathscr E)\le Y_6(\mathscr E)$, and hence
$W_{\mathscr E}(n)\le4{\one}_{\Omega(n)\le2}$.
\end{lemma}

\begin{proof}
Map a contributing four-prefix tuple to its last pair $(p_3,p_4)$.
The quotient roughness forces its first two selected primes to be
the two least factors of $n$. The image lies in the prescribed
rectangle, and removing $p_3p_4$ leaves an $N^{a_{\rm s}}$-rough residual,
so this pair contributes to $Y_6(\mathscr E)$. Since the first two factors
are determined by $n$, two tuples with the same last pair coincide.
The map is therefore injective.
\end{proof}

The fixed prime-factor bound $\lfloor1/{a_{\rm s}}\rfloor=11$ bounds
$|W_{\mathscr E}|$ by $1016$ on rough $n$.  The number of $n\le N$ with a
repeated prime at least $N^{a_{\rm s}}$ is at most $2N^{1-{a_{\rm s}}}$ for large $N$.
After restricting to primes $p<N$ with $(p,N)=1$ and removing the
unit, the squareful and unit costs are at most
$2032N^{1-{a_{\rm s}}}+4=o(H_N)$.  It follows that
\begin{equation}
 4D_{1,2}(N)\ge
       \sum_{\substack{p<N\\p\ {\rm prime}\\(p,N)=1}}W_{\mathscr E}(N-p)
        -o(H_N).
 \label{eq:current-weight-to-count}
\end{equation}

\section{The lower sieve and finite double-sieve transitions}
\label{sec:current-lower}

\subsection{A lower estimate for the original prime sequence}

The ordinary lower trees start at $w={a_{\rm s}}$ and $w={h_{\rm s}}$, with cap
$\nu=249/448$ and fixed small buffers.  Their first exits have even
rank.  The analogue of \eqref{eq:current-upper-exit} is
\begin{equation}
 \sum_d\mu(d){\one}_{d\ {\rm retained}}\xi(d)
 =\prod_p(1-\xi(p))
  -\sum_{e\ {\rm even\ first\ exit}}\xi(e)
                       \prod_{p<p_{\rm last}(e)}(1-\xi(p)).
 \label{eq:current-lower-exit}
\end{equation}
This identity proves both the pointwise lower inequality and the
literal raw-density formula. The selected fourth-rank restorations
preserve even first exits: each cap is chosen before the fourth cube
test $S_3+3t_4<\lambda_4$, and the ordinary residual level exceeds
$p_4^2$. A restored slab and a fifth-rank saving are
assigned disjoint regions in the lower functional.

For a presieve $z_0=N^{\eta_0}$, even exits give the exact projection
direction
\begin{equation}
 V_N(z_0)R_{\rm rough}-V_{\rm full}
  =\sum_{\substack{e\ {\rm full\ even\ exit}\\p_{\rm last}(e)<z_0}}
                              g_N(e)V_N(p_{\rm last}(e))\ge0.
 \label{eq:current-lower-projection}
\end{equation}
Round positive even coefficients inward and negative odd coefficients
outward, retaining the prefix decisions.  Attach beta-minus to positive
coefficients and beta-plus to negative coefficients.  If the rounded
coefficient is $a_{\rm rnd}(d)$, then
$a_{\rm rnd}(d)B_{\rm sign}(n)\le a_{\rm rnd}(d)I_0(n)$ term by term.  Therefore the
physical lower coefficient satisfies
\[
 \sum_{q\mid n}\rho_N^{\rm low}(q)\le
           {\one}_{(n,P_N(N^w))=1},\qquad
 \sum_q\frac{\rho_N^{\rm low}(q)}{\varphi(q)}
              \ge V_{\rm full}-\text{coefficient-realization error}.
\]
Beta-minus need not be nonnegative.  The second inequality uses
\eqref{eq:current-lower-projection} and absolute error control, not
coefficient order alone.  The coefficient theorem supplies
\eqref{eq:current-PP} for this very lower coefficient.

The original prime sequence has a sharp endpoint at $N$. To use the
smooth distribution estimate for a signed lower coefficient, we first
control the endpoint error in absolute value.

\begin{lemma}[Sharp original-prime endpoint]
\label{lem:current-sharp-prime}
Suppose the actual coefficient satisfies $|\rho_N^{\rm low}(q)|\le M$ with
fixed $M$, is supported on squarefree $q\le N^\theta$ with
$\theta<1$ and $(q,N)=1$, and satisfies the permitted
smooth-singleton instance of \eqref{eq:current-PP}.  Then
\begin{equation}
 \sum_q\rho_N^{\rm low}(q)
    \left(\pi(N;q,N)-\frac{\pi(N)}{\varphi(q)}\right)
      =O_A\bigl(N(\log N)^{-A}\bigr).
 \label{eq:current-sharp-pi}
\end{equation}
\end{lemma}

\begin{proof}
Smooth the endpoint over $[N-\delta N,N+\delta N]$, with
$\delta=(\log N)^{-B}$ and $B>A+3$, and apply the smooth singleton
estimate. The absolute boundary and exact-coprime mean estimates in
Appendix~\ref{app:cross-review} bound the difference by
$O(\delta N\log N+\log^2N)$, together with the fixed-power complete
product trim. The equality $p=N$ is impossible for even $N>2$.
These errors are absorbed by the chosen logarithmic accuracy.
\end{proof}

Apply the pointwise lower coefficient to the original prime sequence.
Lemma~\ref{lem:current-sharp-prime} bounds its remainder, while
\eqref{eq:current-raw-density} and the prime number theorem evaluate
its main term. This gives the ordinary lower prime counts used below.

\subsection{Finite defects and actual descendants}

Put
$A_+(s)=sF(s)/(2e^\gamma)$ and
$A_-(s)=sf(s)/(2e^\gamma)$.
For a nonnegative finite marked family $\sigma(B)$ and one common
$Q=N^\lambda$, normalize by
\[
 \Theta_Q=4\operatorname{li}(N)
       \sum_B\frac{\sigma(B)C_{BN}}{\varphi(B)\log(Q/B)}.
\]
Here $C_{BN}$ is the same singular-factor function as in
\eqref{eq:current-normalization}.  A defect $H(s)$ means an upper
bound $(A_+(s)-H(s)+o(1))\Theta_Q$; a defect $h(s)$ means a lower
bound $(A_-(s)+h(s)+o(1))\Theta_Q$.  These are estimates
averaged over the positive marked family, not pointwise in $B$ and
not assertions for a signed outer measure.

For residual exponent $L=\lambda-\log_NB$, write
$\alpha=(1-\log_NB)/L\ge2$.  A singleton child with exponent $Lu$
has
\[
 L'=L(1-u),\qquad \alpha'=1+\frac{\alpha-1}{1-u};
\]
a two-prime child replaces $u$ by $u+v$.  All children retain the
same total $Q$, every full and selected square, and their physical
prime ceiling.  The parameter $\alpha$ increases and is never
reset to $2$ in a switched mass.

Lemma~\ref{sv:finite-seed} proves the initial defect $J_0=1/50$
and the positive transition
$J\mapsto0.020249+0.13J$ on the permitted marked states with
$1\le s\le11/5$. Its hypotheses retain the residual level, the
prefix square conditions and a common comparison cutoff. Four
applications give $J_4=0.023273777353>2327/100000$; this is the seed
for the finite recurrence below.
The finite collection consists of $70$ one-step and $36$ two-step
inequalities, together with the positive Buchstab transitions
\[
 H_{\rm new}(s)=\int_{\max(2,s-1)}^5\frac{h(w)}w\,dw,
 \qquad
 h_{\rm new}(s)=\int_{s-1}^5\frac{H(w)}w\,dw\quad(s\ge2).
\]
The one-step switched cost is $KI/8$, where $I$ is the
nonnegative switched mass of the action specified in
Appendix~\ref{app:lower}.  A two-step action has the
positive defect contribution
\begin{equation}
 \frac15\left[4H(s_0)+H(\phi_1)
  +\sum_{\mathcal D_1}\int\frac{h(\rho)}{u(1-u)}\,du
  +\sum_{\mathcal D_2}\iint\frac{H(\rho)}{uv(1-u-v)}\,du\,dv\right]
 \label{eq:current-w2-feedback}
\end{equation}
in addition to its proved ordinary and positive-switching base.
There are three singleton and four good pair domains.  With
$a_i=1/\phi_i$ and $a_0=1/s_0$, the pair domains are
\[
 \begin{array}{c|c|c}
 &\text{domain}&\rho_i(u,v)\\ \hline
 D_5^+&a_0<u<v<a_2,\ v+2u<1&s_0(1-u-v)\\
 D_6&(a_0,a_1)\mathbin{\times}(a_2,a_3)&s_0(1-u-v)\\
 D_7&a_0<u<v<a_1&(1-u-v)/u\\
 D_8&(a_0,a_1)\mathbin{\times}(a_1,a_2)&(1-u-v)/u.
 \end{array}
\]
The admissibility conditions include $2v<1$ and $v+2u<1$.
The excluded portion of $D_5$ receives only a positive product cost,
not ordinary feedback.  Unit surfaces in the switching terms stay
in their masses.  The envelope used for every actual $\alpha\ge2$ is
$\omega(s)\le\max(1/s,70893/125000)$ for $s\ge1$, with the
conservative extension below $1$ specified in Appendix~\ref{app:lower}.

\paragraph{Two levels of iteration.}
The four scalar steps produce the certified constant used to initialize
the arrays; they precede generation zero of the array recurrence.
The indices in $H_m,h_m$, including $m=248$ and $430$, count only
simultaneous updates of the two cell arrays from that initialization.
Thus the scalar seed iteration and the envelope generations have
different state spaces and iteration counts.

We define the envelopes by simultaneous finite induction. Each
right-hand side uses the preceding generation, takes the maximum of
alternative estimates, and uses a minimum over the full closed support
when an argument crosses a cell boundary. The $200$ closed cells have
endpoints $1+i/50$, $0\le i\le200$. Let $H_m,h_m$ be the step-function
lower envelopes after $m$ generations; at a shared endpoint take the
smaller of the adjacent values. Starting directly from the seed,
$430$ rounds determine all $400$ final entries. These
generation-$430$ arrays are used only in $U_3,U_4$; the remaining
occurrences use generation $248$.

The finite induction also needs a common approximation, not just
finite depth.  Positive transition mass is at most $12$ and the
root coefficient norm at most $7/2$.  A local oriented error
$\delta$ propagates by at most
$\delta(12^{441}-1)/11$ before normalization.  Choose
$\delta$ inside the prescribed final reserve and use one nested
finite refinement for all queries.  Residual exponents after the
finite transitions remain above a fixed positive multiple of
$101^{-440}$, so every positive leaf satisfies the fixed-cutoff
hypothesis of Proposition~\ref{prop:current-positive-product}.
Zero-mass cells are controlled by absolute errors.

\subsection{Evaluation of the eleven-term weight}

The lower functional is
\begin{equation}
 \mathcal C=\frac14
   (3L_1+L_2+L_5+L_6-U_3-U_4-2U_7-U_8-U_9-U_{10}-U_{11}).
 \label{eq:current-root-functional}
\end{equation}
The lower starts are the causal starts at ${a_{\rm s}},{h_{\rm s}}$. The ordinary estimate
at ${h_{\rm s}}$ is inserted after its full least-prime correction is removed
from the reference expression. In $U_3,U_4$ the
$H_{430}$ argument is $\max(2,L/{a_{\rm s}})$ with $L=v(t)-t$ in the enhanced
range. The region $L/{a_{\rm s}}\le2$ has positive measure and uses the value
of the lower envelope at $2$.
The pair lower terms use $h_{248}$ where $L/{a_{\rm s}}\ge2$ and zero below
$2$.  The positive costs $U_7,U_8,U_9$ retain their labelled
multiplicities and are estimated by
Proposition~\ref{prop:current-positive-product}.

For $U_{10},U_{11}$, work on one common positive partition of the
four-prefix domain, with decreasing coordinates $x>y>z>w$.
At each cell use the minimum of the switched kernel
$K\omega((1-S_4)/z)/z$ and the admissible ordinary kernels
$4(1-H_{248}(\max(2,L/w)))/L$. Each ordinary alternative retains
all full and selected square conditions. A refined estimate replaces
its parent estimate, and overlapping alternatives are compared rather
than added. Appendix~\ref{app:lower} gives the states and their common
prefix-preserving approximation.

The restriction in~\eqref{eq:current-paired-weight} requires an estimate
of the lower functional as well as the pointwise injection. On a pair
rectangle $\mathscr B=[t_l,t_r]\times[u_l,u_r]$, put
\[
 R(t,u)=\frac{57}{64}-\frac52t-u,\qquad
 R_{\max}=R(t_l,u_l),\quad s_{\max}=R_{\max}/{a_{\rm s}},\quad
 M_{\mathscr B}=\log(t_r/t_l)\log(u_r/u_l).
\]
Write $\mathcal L_{\mathscr B}$ for the cellwise lower estimator assigned to a
rectangle $\mathscr B$. It vanishes on $\mathscr E_0$. In $\mathscr E_1$ there are
$227$ crossing cells on which it vanishes and $11254$
positive cells with $2\le s_{\min}\le s_{\max}\le9/4$.
Since $\log(s-1)/s$ is increasing there, and $h_{\mathscr B}^{\min}$ denotes the
cellwise $h_{248}$ minimum over the full closed argument support,
the discarded lower estimator, in Chen-quarter units, is
bounded above by
\begin{equation}
 \sum_{\mathscr B\in \mathscr E_1\ {\rm positive}}
       \frac{M_{\mathscr B}}{R_{\max}}\bigl(\log(s_{\max}-1)+h_{\mathscr B}^{\min}\bigr)
       \le 0.002372094.
 \label{eq:current-crop-loss}
\end{equation}

For the negative term, the same selected parents are switched in
the $48$-division four-prefix cover and avoid the refinement on
$40000$ parent cells.
If $A_{\mathscr B}$ is the reference density mass and $s_{\mathscr B}$ its fixed $K$
feedback slope, the cellwise bounds in Appendix~\ref{app:numeric} give
$0\le s_{\mathscr B}\le A_{\mathscr B}/K_0$, $K_0=6.7619$.  Therefore
\[
 A'_{\mathscr B}=A_{\mathscr B}-(K_0-K)s_{\mathscr B}\ge(K/K_0)A_{\mathscr B}.
\]
Dividing by its positive reference mass gives an explicit
pointwise density majorant of the switched kernel.
Restricting this majorant to the complementary positive family
removes at least $KI_{\mathscr E}/8$ in Chen units, where
\begin{equation}
 I_{\mathscr E}=\int_{(t,u)\in \mathscr E}\int_{{a_{\rm s}}<w<z<u}
                      \frac{dt\,du\,dz\,dw}{tu z^2w}.
 \label{eq:current-crop-integral}
\end{equation}
The use of $\omega\ge1/2$ is valid since
$(1-S_4)/z\ge77/38>2$ throughout this family.  For one rectangle
the integral is exactly
\[
 \log(t_r/t_l)\left[
 \frac{\log(u_r/u_l)}{a_{\rm s}}+\frac{\log(u_r/{a_{\rm s}})+2}{u_r}
                         -\frac{\log(u_l/{a_{\rm s}})+2}{u_l}\right].
\]
The finite lower bounds for the gross improvements are
$0.002095161$ on $\mathscr E_0$ and $0.003799514$ on $\mathscr E_1$.
Subtracting \eqref{eq:current-crop-loss}, the additional raw
improvement from $\mathscr E_1$ is
\begin{equation}
          0.003799514-0.002372094
             =0.001427420=\frac{71371}{50000000}.
 \label{eq:current-new-crop-gain}
\end{equation}
The paired restriction removes mass from the positive majorant and
pays for the assigned lower estimators on the same deleted region.
The masks $\mathscr E_0,\mathscr E_1$ are disjoint.

\subsection{A preliminary lower functional}

Denote the current lower endpoint by $R_{\rm base}$ and the bound
for each finite-density and coefficient-realization error by
$B_{\rm base}$; see~\eqref{eq:base-R}--\eqref{eq:base-B}.
There are two further improvements on $\mathscr E$: a larger lower bound
for $\omega$ on an inner fibre, and a smaller payment for the
discarded lower estimator. We treat them separately before combining
their contributions.

\subsection{The paired-domain improvement}
\label{sec:current-refinement}

\begin{lemma}[A global Buchstab lower bound]
\label{lem:current-omega-refinement}
For $s\ge1$ one has $\omega(s)\ge1/2$, and for $s\ge7/3$ one has
$\omega(s)\ge11/20$.
\end{lemma}

\begin{proof}
The bound $1/2$ holds on $[1,2]$ and propagates by
$(s(\omega(s)-1/2))'=\omega(s-1)-1/2$.
For $7/3\le s\le3$, use $\log x\ge2(x-1)/(x+1)$ for $x\ge1$ to get
\[
 \omega(s)=\frac{1+\log(s-1)}s\ge\frac{3s-4}{s^2}\ge\frac{11}{20}.
\]
The last inequality follows because the concave polynomial
$60s-80-11s^2$ is positive at both endpoints.
For $3\le s\le10/3$, the first bound and $\log2\ge2/3$ give
\[
 s\omega(s)\ge1+\log2+(s-3)/2\ge s/2+1/6\ge11s/20.
\]
This establishes the stronger bound on the unit interval
$[7/3,10/3]$. The identity
$(s(\omega(s)-11/20))'=\omega(s-1)-11/20$ then propagates it to
successive unit intervals.
\end{proof}

\paragraph{The extra positive fibre.}
The already selected parent majorant dominates
$K\omega((1-t-u-z-w)/z)/z$ on ${a_{\rm s}}<w<z<u<{h_{\rm s}}<t<{d_{\rm s}}$.
For each rectangle $\mathscr B=[t_l,t_r]\times[u_l,u_r]$ of the same mask,
put
\begin{equation}
 Z_{\mathscr B}=\min\left(u_l,\frac{3(1-t_r-u_r)}{13}\right).
 \label{eq:current-Z}
\end{equation}
The finite exact geometry check gives $Z_{\mathscr B}>{a_{\rm s}}$ for every cell.
Moreover the definition gives
$1-t_r-u_r-2Z_{\mathscr B}\ge(7/3)Z_{\mathscr B}$.  Hence on the inner fibre
${a_{\rm s}}<w<z<Z_{\mathscr B}$ one has $z<u$ and
\[
 1-t-u-z-w\ge1-t_r-u_r-2z\ge\frac73z.
\]
Lemma~\ref{lem:current-omega-refinement} supplies an additional
$1/20$ in $\omega$, beyond the $1/2$ already paid in
\eqref{eq:current-crop-integral}.  The remainder of every fibre
keeps its original bound.  The two nonnegative contributions hold simultaneously.

Integrating the added $1/20$ and including the Chen factor $1/4$
gives at least the following additional coefficient from $\mathscr B$:
\begin{align}
 G_{\mathscr B}
 &=\frac{K}{80}
   \log\frac{t_r}{t_l}\log\frac{u_r}{u_l}
      \int_{a_{\rm s}}^{Z_{\mathscr B}}\frac{\log(z/{a_{\rm s}})}{z^2}\,dz \notag\\
 &=\frac{K}{80}
   \log\frac{t_r}{t_l}\log\frac{u_r}{u_l}
   \left[\frac1{a_{\rm s}}-\frac{1+\log(Z_{\mathscr B}/{a_{\rm s}})}{Z_{\mathscr B}}\right].
 \label{eq:current-extra-fibre}
\end{align}
The inner $w$ integral is $\log(z/{a_{\rm s}})$, and the remaining primitive
is $-(1+\log(z/{a_{\rm s}}))/z$.  Thus there is no permutation factor and
no extra factor $2$.  The fixed rational walls in
\eqref{eq:current-Z} remain within the same admissible finite
positive family.  Restriction uses the positive parent majorant and its assigned cell decomposition.

The directed interval estimate of Appendix~\ref{app:refinement} gives
\begin{equation}
 \sum_{\mathscr B\in \mathscr E_0\cup \mathscr E_1}G_{\mathscr B}>
 \frac{558583}{10^9}=0.000558583=:g.
 \label{eq:current-extra-gross}
\end{equation}
Only this rational lower endpoint is used below.

\paragraph{The discarded positive lower estimator.}
The lower estimator assigned to a rectangle $\mathscr B$ is the maximum of its
constant and chord estimates for the ordinary term, plus its nonnegative
double-sieve defect. The complete sum on $\mathscr E$, including the Chen factor
$1/4$, is denoted by $L_{\mathscr E}$. Its rational fixed-point evaluation gives
\begin{equation}
 L_{\mathscr E}=\frac{5645403422587826010135}{2417851639229258349412352}
 <\frac{466977}{200000000}=0.002334885=:\lambda_{\rm loss}.
 \label{eq:current-literal-loss}
\end{equation}
We next bound this discarded lower functional directly by rational intervals.

To specify the enclosure, on a positive cell use the residual
function $R(t,u)=57/64-(5/2)t-u$ used in the preliminary functional,
and put
\begin{align*}
 r_{\mathscr B}&=R(t_r,u_r),\quad R_{\mathscr B}=R(t_l,u_l),\quad
 s_l=r_{\mathscr B}/{a_{\rm s}},\quad s_h=R_{\mathscr B}/{a_{\rm s}},\\
 f_l&=\frac{\log(s_l-1)}{s_l},\qquad
 f_h=\frac{\log(s_h-1)}{s_h}.
\end{align*}
Here $2\le s_l\le s_h\le9/4$.  Define
\begin{align*}
 M_{\mathscr B}&=\log(t_r/t_l)\log(u_r/u_l),\\
 M_{R,\mathscr B}
 &=\int_{\mathscr B} R(t,u)\frac{dt\,du}{tu}\\
 &=\frac{57}{64}M_{\mathscr B}
   -\frac52(t_r-t_l)\log(u_r/u_l)
   -(u_r-u_l)\log(t_r/t_l).
\end{align*}
The two ordinary alternatives, after division by the Chen
factor $4$, are bounded above by
\begin{equation}
 \frac{f_lM_{\mathscr B}}{{a_{\rm s}}},
 \qquad
 \frac{f_l(R_{\mathscr B}M_{\mathscr B}-M_{R,\mathscr B})+
       f_h(M_{R,\mathscr B}-r_{\mathscr B}M_{\mathscr B})}{a_{\rm s}(R_{\mathscr B}-r_{\mathscr B})}.
 \label{eq:current-floor-proxy}
\end{equation}
Indeed the exact endpoint values $f_l,f_h$ dominate the
lower endpoints used in the functional.  Both coefficient integrals are
nonnegative because $r_{\mathscr B}\le R(t,u)\le R_{\mathscr B}$, and the fixed-point lower roundings can only decrease the estimator value.
The cellwise estimator takes a maximum of its constant and
chord alternatives; taking the maximum of the two upper bounds
in \eqref{eq:current-floor-proxy} preserves the direction.

For the separate defect term, take $h_{\mathscr B}^{\min}$ to be the minimum of
the $h_{248}$ closed-cell values meeting the exact interval
$[s_l,s_h]$, including both adjoining cells at a knot.  This
minimum is at least the one over the outward-rounded lookup
interval.  Therefore $h_{\mathscr B}^{\min}M_{\mathscr B}/R_{\mathscr B}$ is an upper bound for the
cellwise defect estimator, and can be added to the preceding maximum.
All relevant $h_{248}$ entries are nonnegative.  If the outward lookup crosses below $2$ and returns zero, the same
upper bound remains valid.  On the $227$ crossing cells the
constant, eligible chord alternatives and lookup all give zero.
Every $\mathscr E_0$ cell has $s_h\le2$ and likewise has zero cellwise lower estimator.

The proxy sum has a rational upper enclosure strictly below
$\lambda_{\rm loss}$. Appendix~\ref{app:refinement-review} obtains it by range
reduction and alternating logarithm series with explicit remainders.
Thus $\lambda_{\rm loss}$ can replace the payment
in~\eqref{eq:current-crop-loss} without weakening the lower inequality.

The recovered payment is exactly
\begin{equation}
 0.002372094-\lambda_{\rm loss}=\frac{37209}{10^9}=0.000037209.
 \label{eq:current-recovered-loss}
\end{equation}
The retained $L_6$ lower endpoint is the whole $L_6$ lower
endpoint minus $4\lambda_{\rm loss}$, namely
\[
 \frac{1421673188323272382945058469141}
      {472236648286964521369600000000}.
\]
The total removed negative gross in Chen units is now the safe
value $0.002095161+0.003799514+g=0.006453258$.
Combining~\eqref{eq:current-extra-gross} and
\eqref{eq:current-recovered-loss} gives the additional lower coefficient
\begin{equation}
 \Delta=g+(0.002372094-\lambda_{\rm loss})
 =\frac{37237}{62500000}=0.000595792.
 \label{eq:current-refinement-delta}
\end{equation}

We add $\Delta/8=37237/500000000=0.000074474$
to each of the two error allowances.  The two enlarged allowances cover finite density approximation and
coefficient realization, respectively.
The net increment after both extra reserves is
\begin{equation}
 \Delta-2(\Delta/8)=\frac{3\Delta}{4}
                  =\frac{111711}{250000000}=0.000446844.
 \label{eq:current-refinement-net}
\end{equation}
The walls and queries are finite, so their local accuracy
can be chosen inside these positive additional allowances by
the same common approximation and coefficient realization argument used
for the preliminary lower functional.

\subsection{The lower endpoint and the order of limits}
\label{sec:lower-limits}

With the paired-domain improvement, the raw endpoint and each
realization allowance are
\begin{align}
 R_{\rm low}&=R_{\rm base}+\Delta,\label{eq:current-R}\\
 B_{\rm low}&=B_{\rm base}+\Delta/8.\label{eq:current-B}
\end{align}
Their exact values are specified in Appendix~\ref{app:numeric}.
From \eqref{eq:current-weight-to-count}, the actual lower coefficient realization,
finite induction, and evaluation of the paired lower functional,
\begin{equation}
 \liminf_{\substack{N\to\infty\\2\mid N}}
         \frac{D_{1,2}(N)}{H_N}\ge R_{\rm low}-2B_{\rm low}.
 \label{eq:current-lower-derived}
\end{equation}
The two allowances leave
\begin{equation}
                  R_{\rm low}-2B_{\rm low}>c+4.9061\cdot10^{-9},
                  \label{eq:current-after-budgets}
\end{equation}
so in particular
\begin{equation}
                  R_{\rm low}-2B_{\rm low}-c>4.9061\cdot10^{-9}>0.
                  \label{eq:current-lower-slack}
\end{equation}
In particular $D_{1,2}(N)\ge cH_N$ eventually, with $c$ as in~\eqref{eq:current-safe-constants}.

Fix the finite set of estimates and their dependencies, the tail
and compact approximation tolerances, the cutoffs defining fixed compact
parameter domains, the prefix caps, and the support reserves. Next choose the common beta and rough
parameters, fine labels, fixed divisor and moment orders, admissible
$\kappa$, and local logarithmic accuracy. Only then let $N$ tend to
infinity through the even integers. Taking the maximum of the finitely
many thresholds makes all estimates simultaneous. Since $C_N\ge C_2>0$,
fixed power savings and $O_A(N/\log^AN)$ with $A>2$ are uniformly
$o(H_N)$. The finite induction therefore gives the conclusion for every
sufficiently large even target.

\section{Finite numerical estimates}
\label{sec:current-numeric}

The preceding upper and lower constructions reduce the numerical
part of the proof to the following finite inequalities.
Table~\ref{num:certificate-table} gives the correspondence between
these quantities and the certificate fields used to verify them.
Appendix~\ref{app:numeric} defines their integrals, partitions,
directed enclosures, and error allowances.

\begin{lemma}
\label{lem:numerical-endpoints}
The upper density functional has an upper endpoint $E_{\rm up}$ as
in~\eqref{eq:current-E}. It admits finite density and coefficient
realizations with errors at most $u$ each, where $u$ is
as in~\eqref{eq:current-u}, and $E_{\rm up}+2u<K$.
The eleven-term lower functional with paired deletion has a lower
endpoint $R_{\rm low}=R_{\rm base}+\Delta$ as in~\eqref{eq:current-R}.
The corresponding two errors are at most $B_{\rm low}$ each, where $B_{\rm low}$ is
as in~\eqref{eq:current-B}, and $R_{\rm low}-2B_{\rm low}>c$.
The generation-$430$ envelopes used in $U_3,U_4$ are generated by
$430$ simultaneous applications of the positive finite action map
to the stated seed. The generation-$248$ envelopes used in the other
specified terms have the separate enclosures described in
Appendix~\ref{app:numeric}.
\end{lemma}

\begin{proof}
For the upper estimate, evaluate the four terms
\[
 E_{\rm up}=\frac4\Theta+I_1+I_{3,\mathrm{compact}}+I_{3,\mathrm{tail}}
       -\sum_{j=1}^{22}G_j,
\]
using upper endpoints for the four costs and lower endpoints for the
nonnegative savings $G_j$. Each successive change is measured from its
parent; overlapping alternatives are maximized before integration.
The noncompact tail is bounded over its entire range.

For the lower estimate, apply lower endpoints to $L_1,L_2,L_5,L_6$
and upper endpoints to $U_3,U_4,U_7,\ldots,U_{11}$
in~\eqref{eq:current-root-functional}. The changes of lower starts,
restored exit regions, and finite-table corrections are inserted on
their prescribed disjoint regions or as complete parent replacements.
The paired-domain terms are then evaluated by
\eqref{eq:current-crop-integral}, \eqref{eq:current-crop-loss},
\eqref{eq:current-extra-gross}, and~\eqref{eq:current-refinement-delta}.
This gives the rational endpoint $R_{\rm low}$.

All scalar operations are rational. Logarithms are enclosed by series
with explicit remainders; the sieve functions are enclosed by
oriented interpolation estimates. Affine cells are triangulated using
exact vertices and exact polynomial moments. Positive savings may be
bounded on a subset, with the remaining saving discarded; an upper
bound for a negative cost covers the full retained region. The finite
recurrence reads only the preceding generation and rounds every
nonnegative contribution downwards. The detailed enclosure and coverage
statements are given in Appendix~\ref{app:numeric}. Rational subtraction
gives~\eqref{eq:current-upper-slack} and~\eqref{eq:current-lower-slack}.
\end{proof}

\paragraph{Numerical data.}
The finite partitions, recurrence, and directed integral bounds of
Appendix~\ref{app:numeric} are supplied in
\emph{Theorem (1+1.75) on the Goldbach Conjecture --- Numerical
Supplement}, version 20261106~\cite{ChanSupplement}. The appendix describes the data and their certificate
correspondence with Lemma~\ref{lem:numerical-endpoints}.

\section{Ordered bad semiprimes and the final subtraction}
\label{sec:current-transfer}

We now estimate $B_b(N)$. Its dependence on $b$ is contained in
the unsifted mass of the ordered semiprimes; the positive product
estimate then imposes the prime complement.

Fix $1/2<b<1$ and set $\beta=b/(1+b)$.  Products below
$2N^{1-\kappa_0}$ have a negligible count.  For a remaining bad
product $m=rq$, the inequality $r>q^b$ implies $r>m^\beta$;
both primes therefore have a fixed positive exponent.  Use the
coordinates
\[
                       t=\frac{\log r}{\log N},\qquad y=\frac{rq}{N}.
\]
On a fixed radial strip $\eta\le y\le C$, PNT in $q=Ny/r$
followed by reciprocal-prime partial summation in $r$ gives the
unsifted mass element
\begin{equation}
                         \frac{N}{\log N}\,
                              \frac{dt\,dy}{t(1-t)}.
 \label{eq:current-pair-measure}
\end{equation}
The PNT errors are uniform on positive exponent compacta.  A prime
upper bound controls $y<\eta$ by $O_b(\eta N/\log N)$; $N$ tends
to infinity first, then $\eta$ tends to zero.

The exact ordered bad inequalities are
\[
 0<y<1,\qquad
 t\le\frac12+\frac{\log y}{2\log N},\qquad
 t>\beta+\frac{b\log y}{(1+b)\log N}.
\]
Their limiting unsifted mass is
\begin{equation}
 J_b=\int_0^1dy\int_\beta^{1/2}\frac{dt}{t(1-t)}
     =\left[\log\frac{t}{1-t}\right]_\beta^{1/2}
     =\log(1/b).
 \label{eq:current-bad-mass}
\end{equation}
The ordering $r\le q$ counts each composite once.  There is no
factor $2$ and no further singular factor in $J_b$.  The
equal-prime diagonal has size $O(\sqrt N)$ and can be retained or
restored after a squarefree argument.

Smooth the complete-product wall in $y$ and the remaining walls
in $t$.  For any fixed $\epsilon>0$ a positive outer approximation
has limiting mass at most $J_b+\epsilon$ and stays in the admissible
finite family.  A partition of unity repeats no mass.
Proposition~\ref{prop:current-positive-product}, followed by
$N\to\infty$ and then $\epsilon\downarrow0$, gives
\begin{equation}
 \limsup_{\substack{N\to\infty\\2\mid N}}
             \frac{B_b(N)}{H_N}\le K\log(1/b).
 \label{eq:current-bad-bound}
\end{equation}

Unique factorization gives the exact identity
\begin{equation}
                         D_b(N)=D_{1,2}(N)-B_b(N).
 \label{eq:current-exact-subtraction}
\end{equation}
The unit was excluded in \eqref{eq:current-d12}; a convention
including it would change the count by at most one.  Zero is
excluded by $p<N$.  A square is a bad remainder when $b<1$, and
a prime remainder is a good one.  Combining
\eqref{eq:current-lower-derived}, \eqref{eq:current-lower-slack}
and \eqref{eq:current-bad-bound} proves
Theorem~\ref{thm:current-main}.

It remains to verify the strict scalar inequalities in
Corollaries~\ref{cor:three-quarters} and~\ref{cor:current-explicit}. For rational $b$ put $z=(1-b)/(1+b)$.
The positive atanh series and its geometric tail give
\begin{equation}
 2\sum_{j=0}^{95}\frac{z^{2j+1}}{2j+1}
 <\log(1/b)
 <L_{96}(b):=
 2\sum_{j=0}^{95}\frac{z^{2j+1}}{2j+1}
       +\frac{2z^{193}}{193(1-z^2)}.
 \label{eq:current-log-certificate}
\end{equation}
Every endpoint is rational.  Substituting
\eqref{eq:current-safe-constants} gives
\begin{align}
 c-KL_{96}(3/4)&>\frac{16532064}{10^9}>0.0165,\notag\\
 c-KL_{96}(74816964/10^8)
     &>\frac{1001149712517311}{10^{20}}>10^{-5},\notag\\
 c-KL_{96}(74816865/10^8)
     &>\frac{106416285358807}{10^{20}}>10^{-6},
                 \label{eq:current-rational-margins}\\
 c-KL_{96}(9/10)
     &>\frac{124934244400949969098}{10^{20}}>1.24889.\notag
\end{align}
The respective surpluses exceed $3.2064\cdot10^{-5}$,
$1.1497\cdot10^{-8}$, $6.4162\cdot10^{-8}$ and
$4.5244\cdot10^{-4}$.  First choose the fixed realization and
positive-smoothing errors inside these strict margins; then let
the remaining uniform asymptotic errors tend to zero.  Taking the
maximum of the four resulting thresholds proves both corollaries.

\paragraph{Range supplied by the fixed constants.}
For $1/2<b<1$, the right-hand side of \eqref{eq:current-general-theorem}
is positive precisely when $b>\exp(-c/K)$. This value describes the range
supplied by the present pair of coefficient bounds. Extending that range
requires stronger applicable lower or upper estimates, with the
moving-residue and coefficient-class conditions preserved.
